% Source: physical PDF p. 229 (printed p. 206). Coverage: complete Figure
% 28.4 and caption. Uncertainty: none; topology, mass insertion, labels,
% line styles, and arrow directions were reconstructed from the
% full-resolution rendered scan.
\begin{figure}[H]
  \centering
  \begin{tikzpicture}[
    x=1cm,
    y=1cm,
    line cap=round,
    line join=round,
    every node/.style={inner sep=1pt},
    fermion/.style={
      draw,
      line width=0.78pt,
      postaction={decorate},
      decoration={markings,
        mark=at position 0.50 with
          {\arrow{Stealth[length=2.2mm,width=1.6mm]}}}
    },
    scalar/.style={
      draw,
      line width=0.78pt,
      dash pattern=on 3.4pt off 3.4pt
    },
    wavy/.style={
      draw,
      line width=0.72pt,
      decorate,
      decoration={snake,amplitude=0.105cm,segment length=0.31cm}
    }
  ]
    \draw[fermion] (-4.9,0) -- (-2.15,0);
    \draw[scalar] (-2.15,0) -- (2.15,0);
    \draw[fermion] (2.15,0) -- (4.9,0);

    \draw[line width=0.78pt]
      (-2.15,0) .. controls (-1.55,1.62) and (1.55,1.62) .. (2.15,0);
    \draw[wavy]
      (-2.15,0) .. controls (-1.55,1.62) and (1.55,1.62) .. (2.15,0);
    \draw[wavy]
      (0.52,1.18) .. controls (0.92,1.85) and (1.50,2.40) .. (2.30,2.47);

    % Bilinear scalar mixing insertion.
    \draw[line width=1.05pt] (-0.12,-0.17) -- (0.12,0.17);
    \draw[line width=1.05pt] (-0.12,0.17) -- (0.12,-0.17);

    \node[above] at (-3.60,0.04) {$u_L,d_L$};
    \node[above] at (3.60,0.04) {$u_R,d_R$};
    \node[below] at (-1.18,-0.04) {$\mathcal{U},\mathcal{D}$};
    \node[below] at (1.18,-0.04)
      {$\overline{\mathcal{U}}^{\,*},\overline{\mathcal{D}}^{\,*}$};
    \node[anchor=east] at (-2.42,1.23) {gluino};
    \node[anchor=west] at (2.42,0.80) {gluino};
    \node[anchor=west] at (1.62,1.96) {gluon};
  \end{tikzpicture}
  \renewcommand{\thefigure}{28.4}
  \caption{A one-loop diagram for the chromoelectric dipole moment of
  the \(u\) or \(d\) quarks. Here solid lines are quarks; dashed lines
  are squarks; combined solid and wavy lines are gluinos; and the wavy
  line is a gluon. The X represents the insertion of a bilinear
  interaction arising from a trilinear scalar field interaction combined
  with the spontaneous breakdown of
  \(SU(2)\mathbin{\cdot}U(1)\). There are also diagrams in which the
  gluon line is attached to one of the internal squark lines instead of
  to the gluino line.}
  \label{fig:28.4}
\end{figure}
